% =====================================================================
\section{Seminarios prácticos}
% =====================================================================
Esta sección resuelve los catorce problemas de la lista \#2 del tema \cite{gonzalez1}, con su numeración
original (17--30), y las cuestiones del parcial de 2025/26 \cite{parcial} que corresponden a la primera
ley. Cada problema sigue la misma estructura que en el Tema 1: el \emph{enunciado}, la \emph{idea} que guía la
resolución, la \emph{resolución} paso a paso con todo el desarrollo simbólico y, cuando aporta, una
\emph{observación} y una figura. Salvo indicación contraria, $R=\SI{8,314}{J/(mol.K)}$,
$\SI{1}{bar}=\SI{e5}{Pa}$ y $\SI{1}{atm}=\SI{101325}{Pa}$.

\subsection{Cómo plantear un problema del Tema 2}

Los problemas del tema se reducen a cinco tareas (tabla~\ref{tab:metodos}). La primera decisión es si el
proceso es \emph{reversible}, porque eso decide qué presión entra en el trabajo.

\begin{table}[htbp]
\centering
\caption{Mapa de decisión: qué método conviene según lo que pide el enunciado.}
\label{tab:metodos}
\small
\begin{tabularx}{\textwidth}{Xlr}
\toprule
Lo que da / pide el enunciado & Método & Problemas\\
\midrule
Proceso reversible con $p(V)$ conocida; hallar $W$ & M1. Trabajo reversible & 17--19, 21, 23\\
Expansión o compresión contra $p_{\mathrm{ext}}$ fija, por etapas & M2. Trabajo irreversible & 22, 27, 29\\
Dos de $\Delta U$, $Q$, $W$ conocidos; hallar el tercero & M3. Balance de energía & 21--23, 29, P.3, P.4\\
$U$ y ecuación de estado dadas; hallar $C_V$, $C_p$, $\gamma$, $Q$ & M4. Calorimetría & 24, 25, 28(a), 30, P.2\\
Adiabática reversible; hallar $T$, $p$, $V$ o $n$ & M5. Adiabáticas & 20, 26, 28(b)\\
\bottomrule
\end{tabularx}
\end{table}

\begin{recuadro}{Los cinco métodos}
\begin{enumerate}[label=\textbf{\color{navy}M\arabic*.},leftmargin=2.2em,itemsep=4pt]
\item \emph{Trabajo reversible.} Escribir $p$ como función de $V$ a lo largo del proceso y calcular
  $W=-\int_{V_1}^{V_2}p(V)\,\dd V$ (tabla~\ref{tab:trabajos}). En un gas ideal isotermo, $p=nRT/V$.
\item \emph{Trabajo irreversible.} En cada etapa, $W=-p_{\mathrm{ext}}(V_f-V_i)$ con la presión
  \emph{exterior}, no la del gas; los volúmenes extremos salen de la ecuación de estado.
\item \emph{Balance de energía.} Fijar signos (entra $>0$) y aplicar $\Delta U=Q+W$. Si el gas es ideal,
  $\Delta U=nC_V\Delta T$ en cualquier proceso; si hay agitador, $W=W_{\mathrm{conf}}+W_{\mathrm{dis}}$.
\item \emph{Calorimetría.} $C_V=\pdi UTV$; $C_p=C_V+l\,\pdi VTp$ con $l=\pdi UVT+p$; $Q=\int C\,\dd T$
  a lo largo del proceso correspondiente.
\item \emph{Adiabáticas.} Gas perfecto: ecuaciones de Poisson \eqref{eq:poisson2} y $W=\Delta U$. Gas
  cualquiera: $C_V\dd T+l\,\dd V=0$, separar variables e integrar.
\end{enumerate}
\end{recuadro}

\subsection{El balance de energía paso a paso}\label{sec:receta}

Casi todos los problemas del tema son un balance de energía. Escribir explícitamente la secuencia
siguiente es lo que demuestra en el examen que se sabe resolver el problema.

\paso{Paso 0: sistema y frontera} Decir qué es el sistema (el gas, el gas más el agitador\dots) y qué
paredes tiene: rígida o móvil, adiabática o diatérmica.

\paso{Paso 1: proceso} Clasificarlo: ¿reversible o no? ¿Isotermo, isóbaro, isócoro, adiabático,
politrópico, contra $p_{\mathrm{ext}}$ fija, cíclico? Cada restricción elimina una incógnita
(tabla~\ref{tab:procesos}).

\paso{Paso 2: estados} Completar $(p,V,T)$ de los estados inicial y final con la ecuación de estado y,
si el proceso es adiabático reversible, con Poisson. Pasar todo a unidades SI.

\paso{Paso 3: trabajo} Reversible: $W=-\int p\,\dd V$ con la $p(V)$ del proceso (M1). Irreversible:
$W=-\int p_{\mathrm{ext}}\,\dd V$ (M2). Añadir el trabajo disipativo si lo hay.

\paso{Paso 4: energía interna} $\Delta U$ sale de la ecuación energética y solo depende de los estados:
para un gas perfecto, $\Delta U=nC_V\Delta T$, aunque el proceso no sea isócoro.

\paso{Paso 5: balance} Despejar la magnitud que falta de $\Delta U=Q+W$.

\paso{Paso 6: comprobación} Signos: expansión, $W<0$; calor cedido, $Q<0$. Si el mismo cambio de estado
se recorre por dos caminos, $\Delta U$ debe coincidir aunque $Q$ y $W$ no lo hagan.

\begin{figure}[htbp]
\centering
\begin{tikzpicture}[font=\footnotesize,
   p/.style={draw=navy,fill=fondo,align=center,rounded corners=1pt,inner sep=3pt,text width=2.55cm,minimum height=1.0cm},
   d/.style={draw=acento,fill=white,diamond,aspect=2.4,align=center,inner sep=1pt,font=\scriptsize},
   f/.style={-{Stealth[length=4pt]},navy}]
  \node[p] (s0) at (0,0) {\textbf{0} Sistema\\ y paredes};
  \node[p] (s1) at (3.15,0) {\textbf{1} Proceso\\ y restricciones};
  \node[p] (s2) at (6.3,0) {\textbf{2} Estados\\ $(p,V,T)$ en SI};
  \node[d] (dd) at (9.6,0) {¿reversible?};
  \node[p,text width=3.0cm] (sA) at (12.55,0.95) {\textbf{3A} Reversible, M1\\ $W=-\!\int p\,\dd V$};
  \node[p,text width=3.0cm] (sB) at (12.55,-0.95) {\textbf{3B} Irreversible, M2\\ $W=-\!\int p_{\mathrm{ext}}\dd V$};
  \node[p] (s4) at (9.6,-2.6) {\textbf{4} $\Delta U$ de la\\ ec. energética};
  \node[p] (s5) at (6.3,-2.6) {\textbf{5} Balance\\ $\Delta U=Q+W$};
  \node[p] (s6) at (3.15,-2.6) {\textbf{6} Signos y\\ comprobación};
  \draw[f] (s0) -- (s1); \draw[f] (s1) -- (s2); \draw[f] (s2) -- (dd);
  \draw[f] (dd.north) |- node[pos=0.25,left=2pt,font=\scriptsize\bfseries,acento] {sí} (sA.west);
  \draw[f] (dd.south) |- node[pos=0.25,left=2pt,font=\scriptsize\bfseries,acento] {no} (sB.west);
  \draw[f] (sA.east) -- ++(0.4,0) |- (s4.east);
  \draw[navy] (sB.east) -- ++(0.4,0);
  \draw[f] (s4) -- (s5); \draw[f] (s5) -- (s6);
\end{tikzpicture}
\caption{Esquema del balance de energía. La única bifurcación es la reversibilidad, que decide si el trabajo
se calcula con la presión del gas o con la exterior; $\Delta U$ no depende de ella.}
\label{fig:receta}
\end{figure}

\begin{observacion}[El error más frecuente]
Usar $-\int p\,\dd V$ con la presión del gas en un proceso irreversible. Si el gas se expande contra una
presión exterior fija, solo $p_{\mathrm{ext}}$ está definida durante el proceso; la presión del gas solo se
conoce en los estados de equilibrio inicial y final (problemas 27 y 29).
\end{observacion}

% ---------------------------------------------------------------------
\subsection{Problema 17. Compresión con \texorpdfstring{$p=a+bV$}{p = a + bV}}
% ---------------------------------------------------------------------
\begin{problema}{17}{Proceso lineal en el plano $p$--$V$}
En un cilindro--pistón se comprime reversiblemente un gas desde \SI{0,5}{MPa} y \SI{25e-3}{m^3} hasta
\SI{2,0}{MPa} y \SI{5e-3}{m^3}, según $p=a+bV$. (a) Calcúlense $a$ y $b$. (b) Calcúlese el trabajo, en julios.
\end{problema}

\begin{resolucion}
\idea Dos estados fijan la recta; el trabajo es método M1 con $p(V)$ lineal.

\paso{Paso 1: constantes} Imponiendo la recta en los dos estados,
\[
  \left.\begin{aligned}0{,}5&=a+25\cdot10^{-3}\,b\\ 2{,}0&=a+5\cdot10^{-3}\,b\end{aligned}\right\}
  \;\Longrightarrow\;
  b=\frac{2{,}0-0{,}5}{(5-25)\cdot10^{-3}}=\frac{1{,}5}{-0{,}02}=\boxed{-\SI{75}{MPa/m^3}},
\]
y sustituyendo en la primera, $a=0{,}5-25\cdot10^{-3}\cdot(-75)=0{,}5+1{,}875=\boxed{\SI{2,375}{MPa}}$.

\paso{Paso 2: trabajo} Integrando la recta entre $V_1=\SI{0,025}{m^3}$ y $V_2=\SI{0,005}{m^3}$:
\[
  W=-\int_{V_1}^{V_2}(a+bV)\,\dd V=-\Bigl[a(V_2-V_1)+\frac b2\bigl(V_2^2-V_1^2\bigr)\Bigr].
\]
Con $V_2-V_1=-0{,}02$ y $V_2^2-V_1^2=(25-625)\cdot10^{-6}=-6\cdot10^{-4}\ \mathrm{m^6}$:
\[
  W=-\bigl[2{,}375\cdot10^{6}\cdot(-0{,}02)+(-37{,}5\cdot10^{6})(-6\cdot10^{-4})\bigr]
   =-\bigl[-47\,500+22\,500\bigr]=\boxed{\SI{25}{kJ}}.
\]

\paso{Paso 3: comprobación} El área bajo una recta es la de un trapecio:
$|W|=\tfrac12(p_1+p_2)|V_2-V_1|=\tfrac12(2{,}5\cdot10^{6})(0{,}02)=\SI{25}{kJ}$. El signo es positivo porque
es una compresión: el entorno realiza trabajo sobre el gas.
\end{resolucion}

\begin{figure}[htbp]
\centering
\begin{tikzpicture}
\begin{axis}[graf,width=7.6cm,height=5.4cm,xmin=0,xmax=30,ymin=0,ymax=2.6,
    xtick={5,25},ytick={0.5,2},yticklabels={0{,}5,2{,}0},
    xlabel={$V$ (L)},ylabel={$p$ (MPa)}]
  \addplot[draw=none,fill=fondomedio,domain=5:25,samples=2] {2.375-0.075*x} \closedcycle;
  \addplot[draw=none,pattern=north east lines,pattern color=acentoclaro,domain=5:25,samples=2] {2.375-0.075*x} \closedcycle;
  \addplot[navy,domain=5:25,samples=2] {2.375-0.075*x};
  \draw[navy,-{Stealth[length=5pt]}] (axis cs:19,{2.375-0.075*19+0.17}) -- (axis cs:13,{2.375-0.075*13+0.17});
  \fill[rojo] (axis cs:25,0.5) circle (1.8pt) node[right=2pt,font=\scriptsize,black] {1};
  \fill[rojo] (axis cs:5,2) circle (1.8pt) node[above right=1pt,font=\scriptsize,black] {2};
  \node[font=\small,navy] at (axis cs:15,0.55) {$W=+\SI{25}{kJ}$};
\end{axis}
\end{tikzpicture}
\caption{Problema 17. Compresión a lo largo de la recta $p=a+bV$; el trabajo es el área del trapecio, positiva
porque el volumen disminuye.}
\label{fig:p17}
\end{figure}

% ---------------------------------------------------------------------
\subsection{Problema 18. Expansiones politrópicas}
% ---------------------------------------------------------------------
\begin{problema}{18}{Tres exponentes}
Un gas se expande reversiblemente desde \SI{3}{bar} y \SI{0,1}{m^3} hasta \SI{0,2}{m^3}, con $pV^a=\cte$.
Calcúlese, en kJ, el trabajo para (a) $a=0$, (b) $a=1$ y (c) $a=1{,}5$.
\end{problema}

\begin{resolucion}
\idea Método M1 con las tres filas de la tabla~\ref{tab:trabajos}. Datos: $p_1=\SI{3e5}{Pa}$,
$V_1=\SI{0,1}{m^3}$, $V_2=\SI{0,2}{m^3}$, luego $p_1V_1=\SI{30}{kJ}$ y $V_2/V_1=2$.

\paso{(a) $a=0$, isóbaro}
$W=-p_1(V_2-V_1)=-3\cdot10^{5}\cdot0{,}1=\boxed{-\SI{30}{kJ}}$.

\paso{(b) $a=1$, $pV=\cte$} Con $p=p_1V_1/V$,
\[
  W=-p_1V_1\int_{V_1}^{V_2}\frac{\dd V}{V}=-p_1V_1\ln\frac{V_2}{V_1}=-30\ln2\ \mathrm{kJ}=\boxed{-\SI{20,79}{kJ}}.
\]

\paso{(c) $a=1{,}5$} Primero la presión final, $p_2=p_1(V_1/V_2)^{1{,}5}=3\cdot2^{-1{,}5}=\SI{1,0607}{bar}$, de
modo que $p_2V_2=\SI{21,213}{kJ}$. Con la fórmula del politrópico,
\[
  W=\frac{p_2V_2-p_1V_1}{a-1}=\frac{21{,}213-30}{0{,}5}\ \mathrm{kJ}=\boxed{-\SI{17,57}{kJ}}.
\]

\paso{Comprobación} Cuanto mayor es $a$, más deprisa cae $p$ y menor es el área bajo la curva:
$|W_a|>|W_b|>|W_c|$ (figura~\ref{fig:p18}). Los tres son negativos: el gas entrega trabajo.
\end{resolucion}

\begin{figure}[htbp]
\centering
\begin{tikzpicture}
\begin{axis}[graf,width=7.6cm,height=5.4cm,xmin=0.08,xmax=0.22,ymin=0,ymax=3.6,
    xtick={0.1,0.15,0.2},xticklabels={0{,}10,0{,}15,0{,}20},ytick={1.0607,1.5,3},yticklabels={1{,}06,1{,}5,3},
    xlabel={$V$ (m$^3$)},ylabel={$p$ (bar)},leyenda abajo]
  \addplot[navy,domain=0.1:0.2,samples=2] {3};
  \addplot[acento,densely dashed,domain=0.1:0.2,samples=60] {0.3/x};
  \addplot[rojo,domain=0.1:0.2,samples=60] {3*(0.1/x)^1.5};
  \fill[black] (axis cs:0.1,3) circle (1.6pt);
  \legend{(a) $a=0$,(b) $a=1$,(c) $a=1{,}5$}
\end{axis}
\end{tikzpicture}
\caption{Problema 18. Las tres expansiones parten del mismo estado; el trabajo entregado es el área bajo cada
curva entre \SI{0,1}{m^3} y \SI{0,2}{m^3}.}
\label{fig:p18}
\end{figure}

% ---------------------------------------------------------------------
\subsection{Problema 19. Expansión isoterma de gases reales}
% ---------------------------------------------------------------------
\begin{problema}{19}{Trabajo isotermo con tres ecuaciones de estado}
Calcúlese el trabajo realizado por $n$ moles en una expansión isoterma reversible de $V_0$ a $2V_0$ si su
ecuación de estado es (1) $p(V-nb)=nRT$; (2) $pV=nRT\bigl(1+\frac{nB}{V}\bigr)$; (3)
$\bigl(p+\frac{n^2a}{V^2}\bigr)(V-nb)=nRT$.
\end{problema}

\begin{resolucion}
\idea Método M1: despejar $p(V)$ a $T$ fija e integrar término a término. Primitivas que aparecen:
$\int\frac{\dd V}{V-c}=\ln(V-c)$ y $\int V^{-2}\dd V=-V^{-1}$.

\paso{(1) Clausius} $p=\dfrac{nRT}{V-nb}$:
\[
  W=-nRT\int_{V_0}^{2V_0}\frac{\dd V}{V-nb}=\boxed{-nRT\ln\frac{2V_0-nb}{V_0-nb}}.
\]

\paso{(2) Virial} $p=\dfrac{nRT}{V}+\dfrac{n^2BRT}{V^2}$:
\[
  W=-nRT\ln\frac{2V_0}{V_0}-n^2BRT\Bigl[-\frac1V\Bigr]_{V_0}^{2V_0}
   =-nRT\ln2-n^2BRT\Bigl(\frac1{V_0}-\frac1{2V_0}\Bigr),
\]
es decir,
\[
  \boxed{W=-nRT\Bigl(\ln2+\frac{nB}{2V_0}\Bigr)}.
\]

\paso{(3) Van der Waals} $p=\dfrac{nRT}{V-nb}-\dfrac{n^2a}{V^2}$: el primer término es el de (1) y el segundo
cambia de signo al integrar,
\[
  W=-nRT\ln\frac{2V_0-nb}{V_0-nb}+n^2a\Bigl(\frac1{V_0}-\frac1{2V_0}\Bigr)
   =\boxed{-nRT\ln\frac{2V_0-nb}{V_0-nb}+\frac{n^2a}{2V_0}}.
\]

\paso{Comprobación} Con $b=B=a=0$ los tres dan $-nRT\ln2$, el gas ideal. El término $+n^2a/(2V_0)$ reduce
el trabajo entregado: las atracciones moleculares disminuyen la presión del gas.
\end{resolucion}

% ---------------------------------------------------------------------
\subsection{Problema 20. Expansión adiabática del argón}
% ---------------------------------------------------------------------
\begin{problema}{20}{Moles a partir del trabajo adiabático}
Cierta cantidad de Ar, gas perfecto, se expande reversible y adiabáticamente desde \SI{37}{\celsius} hasta
triplicar su volumen, produciendo un trabajo de \SI{3600}{J}. ¿Cuántos moles se expanden?
\end{problema}

\begin{resolucion}
\idea Método M5: Poisson da $T_2$; el balance con $Q=0$ liga el trabajo con $n$.

\paso{Paso 0: datos} Ar monoatómico: $\bar C_V=\tfrac32R$, $\gamma=5/3$. $T_1=\SI{310,15}{K}$, $V_2=3V_1$. El gas
\emph{produce} trabajo, luego con la convención del curso $W=-\SI{3600}{J}$.

\paso{Paso 1: temperatura final} De \eqref{eq:poisson2} con $\gamma-1=2/3$:
\[
  T_2=T_1\Bigl(\frac{V_1}{V_2}\Bigr)^{\gamma-1}=310{,}15\cdot3^{-2/3}\ \mathrm{K}
     =\frac{310{,}15}{2{,}0801}\ \mathrm{K}=\SI{149,10}{K}.
\]

\paso{Paso 2: balance} Con $Q=0$, $W=\Delta U=n\bar C_V(T_2-T_1)$, y despejando $n$:
\[
  n=\frac{W}{\bar C_V(T_2-T_1)}=\frac{-3600}{\tfrac32\cdot8{,}314\cdot(149{,}10-310{,}15)}\ \mathrm{mol}
   =\frac{3600}{12{,}471\cdot161{,}05}\ \mathrm{mol}=\boxed{\SI{1,793}{mol}}.
\]

\paso{Comprobación} $\Delta U<0$ y $T_2<T_1$: el gas se enfría porque el trabajo que entrega sale de su
energía interna.
\end{resolucion}

% ---------------------------------------------------------------------
\subsection{Problema 21. Proceso \texorpdfstring{$pV^2=\cte$}{pV2 = cte} con calor}
% ---------------------------------------------------------------------
\begin{problema}{21}{Energía interna a partir de $Q$ y $W$}
Un gas a \SI{6}{bar} y \SI{0,05}{m^3} sigue un proceso reversible con $pV^2=\cte$ hasta \SI{1,5}{bar}. Calcúlese
$\Delta U$ si recibe \SI{5}{kJ} de calor.
\end{problema}

\begin{resolucion}
\idea Métodos M1 (trabajo del politrópico) y M3 (balance).

\paso{Paso 1: volumen final} De $p_1V_1^2=p_2V_2^2$:
$V_2=V_1\sqrt{p_1/p_2}=0{,}05\sqrt{4}=\SI{0,1}{m^3}$.

\paso{Paso 2: trabajo} Con $m=2$: $p_1V_1=\SI{30}{kJ}$ y $p_2V_2=1{,}5\cdot10^{5}\cdot0{,}1=\SI{15}{kJ}$, luego
\[
  W=\frac{p_2V_2-p_1V_1}{m-1}=\frac{15-30}{1}\ \mathrm{kJ}=-\SI{15}{kJ}.
\]
Directamente: $W=-p_1V_1^2\int_{V_1}^{V_2}V^{-2}\dd V=-p_1V_1^2\bigl(\tfrac1{V_1}-\tfrac1{V_2}\bigr)
=-1500\cdot(20-10)\ \mathrm J=-\SI{15}{kJ}$.

\paso{Paso 3: balance}
$\Delta U=Q+W=5-15\ \mathrm{kJ}=\boxed{-\SI{10}{kJ}}$. El gas entrega más trabajo del calor que recibe.
\end{resolucion}

% ---------------------------------------------------------------------
\subsection{Problema 22. Pistón con agitador}
% ---------------------------------------------------------------------
\begin{problema}{22}{Trabajo de configuración y disipativo}
Un cilindro vertical contiene un gas comprimido por un pistón sin rozamiento de \SI{3000}{N} de peso y
\SI{50e-4}{m^2} de área; fuera, $p_{\mathrm{atm}}=\SI{0,95}{bar}$. Un agitador comunica al gas \SI{6800}{J}. Si el gas
cede \SI{8,7}{kJ} de calor y $\Delta U=-\SI{1}{kJ}$, ¿cuánto se mueve el pistón?
\end{problema}

\begin{resolucion}
\idea Método M3 con $W=W_{\mathrm{conf}}+W_{\mathrm{dis}}$; el trabajo de configuración es contra la presión
exterior constante (M2), que incluye el peso del pistón.

\paso{Paso 0: sistema y signos} Sistema: el gas. $W_{\mathrm{dis}}=+\SI{6800}{J}$ (el agitador da trabajo),
$Q=-\SI{8700}{J}$ (cede calor), $\Delta U=-\SI{1000}{J}$.

\paso{Paso 1: trabajo total} Del balance,
$W=\Delta U-Q=-1000+8700\ \mathrm J=\SI{7700}{J}$.

\paso{Paso 2: trabajo de configuración}
$W_{\mathrm{conf}}=W-W_{\mathrm{dis}}=7700-6800\ \mathrm J=\SI{900}{J}>0$: el pistón baja.

\paso{Paso 3: presión exterior} El pistón está en equilibrio bajo la atmósfera y su peso:
\[
  p_{\mathrm{ext}}=p_{\mathrm{atm}}+\frac{mg}{A}=95\,000+\frac{3000}{0{,}005}\ \mathrm{Pa}=\SI{695000}{Pa}.
\]

\paso{Paso 4: desplazamiento} Con $W_{\mathrm{conf}}=-p_{\mathrm{ext}}\Delta V$ y $\Delta V=A\,\Delta z$:
\[
  \Delta z=-\frac{W_{\mathrm{conf}}}{p_{\mathrm{ext}}A}=-\frac{900}{695\,000\cdot0{,}005}\ \mathrm m
         =\boxed{-\SI{0,259}{m}},
\]
es decir, el pistón desciende \SI{25,9}{cm}.
\end{resolucion}

\begin{observacion}[Por qué baja]
El agitador aporta \SI{6,8}{kJ} pero el gas cede \SI{8,7}{kJ} de calor: si $\Delta U$ solo baja \SI{1}{kJ}, el entorno
tiene que aportar los \SI{0,9}{kJ} restantes comprimiendo el gas.
\end{observacion}

% ---------------------------------------------------------------------
\subsection{Problema 23. Tres caminos entre dos estados}
% ---------------------------------------------------------------------
\begin{problema}{23}{Isoterma, recta y quebrada}
Dos moles de O$_2$, gas perfecto, pasan reversiblemente de presión $p$ a $3p$, ambos a la misma temperatura
$T$, por (a) una isoterma, (b) una recta en el diagrama $pV$ y (c) un isócoro seguido de un isóbaro.
Calcúlense, en función de $T$, los trabajos y los calores.
\end{problema}

\begin{resolucion}
\idea Los tres caminos tienen los mismos extremos y $T_1=T_2$, luego $\Delta U=0$ en los tres (gas ideal) y
$Q=-W$. Solo hay que calcular tres áreas (método M1).

\paso{Paso 0: estados} $V_1=\dfrac{2RT}{p}$, y como la temperatura es la misma, $V_2=\dfrac{2RT}{3p}=\dfrac{V_1}3$.
Además $pV_1=2RT$.

\paso{(a) Isoterma}
\[
  W_a=-nRT\ln\frac{V_2}{V_1}=-2RT\ln\frac13=\boxed{2RT\ln3\simeq2{,}20\,RT},\qquad Q_a=-W_a .
\]

\paso{(b) Recta} El área bajo una recta es un trapecio de alturas $p$ y $3p$:
\[
  W_b=-\frac{p+3p}{2}\,(V_2-V_1)=-2p\Bigl(-\frac{2V_1}{3}\Bigr)=\frac43\,pV_1=\boxed{\frac83\,RT\simeq2{,}67\,RT},\qquad Q_b=-W_b .
\]

\paso{(c) Isócoro + isóbaro} El isócoro ($V_1$, de $p$ a $3p$) no hace trabajo; el isóbaro a $3p$ de $V_1$ a
$V_1/3$ sí:
\[
  W_c=0-3p\Bigl(\frac{V_1}{3}-V_1\Bigr)=2pV_1=\boxed{4RT},\qquad Q_c=-W_c .
\]

\paso{Comprobación} $W_a<W_b<W_c$: cuanto más alta va la curva, mayor es el área (figura~\ref{fig:p23}).
$\Delta U=0$ en los tres, pero $Q$ y $W$ dependen del camino.
\end{resolucion}

\begin{figure}[htbp]
\centering
\begin{tikzpicture}
\begin{axis}[graf,width=7.4cm,height=5.6cm,xmin=0.2,xmax=1.15,ymin=0.5,ymax=3.5,
    xtick={0.3333,1},xticklabels={$V_1/3$,$V_1$},ytick={1,3},yticklabels={$p$,$3p$},
    xlabel={$V$},ylabel={$p$},leyenda fuera]
  \addplot[navy,domain=0.3333:1,samples=60] {1/x};
  \addplot[acento,densely dashed,domain=0.3333:1,samples=2] {1+3*(1-x)};
  \addplot[rojo,densely dotted,line width=1.1pt] coordinates {(1,1) (1,3) (0.3333,3)};
  \fill[black] (axis cs:1,1) circle (1.7pt) node[below right=1pt,font=\scriptsize] {1};
  \fill[black] (axis cs:0.3333,3) circle (1.7pt) node[above left=1pt,font=\scriptsize] {2};
  \legend{(a) isoterma,(b) recta,(c) isócoro + isóbaro}
\end{axis}
\end{tikzpicture}
\caption{Problema 23. Tres caminos reversibles entre los mismos estados. El trabajo recibido es el área bajo
cada uno: $W_a=2{,}20\,RT<W_b=2{,}67\,RT<W_c=4\,RT$.}
\label{fig:p23}
\end{figure}

% ---------------------------------------------------------------------
\subsection{Problema 24. Calor con \texorpdfstring{$\bar C_p$}{Cp} variable}
% ---------------------------------------------------------------------
\begin{problema}{24}{Integral de la capacidad calorífica}
El calor molar a presión constante de un gas es $\bar C_p=a+bT-c/T^2$. ¿Qué calor recibe el gas en un proceso
isóbaro reversible en el que \SI{2}{mol} pasan de $T$ a $2T$?
\end{problema}

\begin{resolucion}
\idea Método M4: a presión constante $\dbar Q=n\bar C_p\,\dd T$ (definición \eqref{eq:CVCp}), y se integra
término a término.

\paso{Paso 1: integral} Con primitivas $aT'$, $\tfrac12bT'^2$ y $+c/T'$ (porque $\int T'^{-2}\dd T'=-T'^{-1}$):
\[
  Q=n\int_{T}^{2T}\Bigl(a+bT'-\frac{c}{T'^2}\Bigr)\dd T'
   =n\Bigl[aT'+\frac b2T'^2+\frac c{T'}\Bigr]_{T}^{2T}.
\]

\paso{Paso 2: evaluación} Con $n=2$:
\[
  Q=2\Bigl[a(2T-T)+\frac b2\bigl(4T^2-T^2\bigr)+c\Bigl(\frac1{2T}-\frac1T\Bigr)\Bigr]
   =2\Bigl[aT+\frac32bT^2-\frac{c}{2T}\Bigr]=\boxed{2aT+3bT^2-\frac cT}.
\]
\end{resolucion}

% ---------------------------------------------------------------------
\subsection{Problema 25. Capacidades del gas de Van der Waals}
% ---------------------------------------------------------------------
\begin{problema}{25}{$\bar C_V$ y $\bar C_p$ con ecuación energética}
Un gas obedece la ecuación de Van der Waals, $(p+a/\Vm^2)(\Vm-b)=RT$, y su ecuación energética es
$\bar U=cT-a/\Vm$, con $c$ constante. Calcúlense $\bar C_V$ y $\bar C_p$.
\end{problema}

\begin{resolucion}
\idea Método M4: $\bar C_V$ directo de $\bar U$; $\bar C_p=\bar C_V+l\,\pdi{\Vm}{T}{p}$ con
$l=\pdi{\bar U}{\Vm}{T}+p$ \eqref{eq:CV}.

\paso{Paso 1: $\bar C_V$} $\bar C_V=\pdi{\bar U}{T}{\Vm}=\boxed{c}$, constante.

\paso{Paso 2: coeficiente $l$} $\pdi{\bar U}{\Vm}{T}=\dfrac{a}{\Vm^2}$, y por la ecuación de estado
\[
  l=\frac{a}{\Vm^2}+p=\frac{RT}{\Vm-b}.
\]

\paso{Paso 3: derivada $\pdi{\Vm}{T}{p}$} No se despeja $\Vm$ (cúbica); se despeja $T$ y se usa la recíproca
(Tema 1). De $T=\frac1R(p+a/\Vm^2)(\Vm-b)$, derivando a $p$ fija con la regla del producto,
\[
  \pd{T}{\Vm}{p}=\frac1R\Bigl[\Bigl(p+\frac a{\Vm^2}\Bigr)-\frac{2a}{\Vm^3}(\Vm-b)\Bigr]
   =\frac1R\Bigl[\frac{RT}{\Vm-b}-\frac{2a(\Vm-b)}{\Vm^3}\Bigr]
   =\frac{RT\Vm^3-2a(\Vm-b)^2}{R\Vm^3(\Vm-b)},
\]
luego $\pdi{\Vm}{T}{p}=\dfrac{R\Vm^3(\Vm-b)}{RT\Vm^3-2a(\Vm-b)^2}$.

\paso{Paso 4: $\bar C_p$} El factor $\Vm-b$ se simplifica:
\[
  \bar C_p=c+\frac{RT}{\Vm-b}\cdot\frac{R\Vm^3(\Vm-b)}{RT\Vm^3-2a(\Vm-b)^2}
  =\boxed{c+\frac{R^2T\Vm^3}{RT\Vm^3-2a(\Vm-b)^2}}.
\]

\paso{Comprobación} Con $a=0$: $\bar C_p=c+R$, la relación de Mayer. Para $a>0$, $\bar C_p-\bar C_V>R$.
\end{resolucion}

% ---------------------------------------------------------------------
\subsection{Problema 26. Adiabáticas de un gas de Clausius}
% ---------------------------------------------------------------------
\begin{problema}{26}{Adiabática con ecuaciones no ideales}
Hállese la ecuación de las adiabáticas en un diagrama $p\Vm$ de un gas con $p(\Vm-b)=AT$ y $\bar U=CT+D$,
siendo $b$, $A$, $C$ y $D$ constantes positivas.
\end{problema}

\begin{resolucion}
\idea Método M5 general \eqref{eq:adgeneral}: $\bar C_V\dd T+l\,\dd\Vm=0$, separar variables e integrar.

\paso{Paso 1: coeficientes} $\bar C_V=\pdi{\bar U}{T}{\Vm}=C$ y, como $\bar U$ no depende de $\Vm$,
$l=0+p=\dfrac{AT}{\Vm-b}$.

\paso{Paso 2: separación} Dividiendo $C\,\dd T+\dfrac{AT}{\Vm-b}\,\dd\Vm=0$ entre $CT$:
\[
  \frac{\dd T}{T}+\frac AC\,\frac{\dd\Vm}{\Vm-b}=0 .
\]

\paso{Paso 3: integración}
$\ln T+\frac AC\ln(\Vm-b)=\cte$, es decir $T(\Vm-b)^{A/C}=\cte$.

\paso{Paso 4: paso al plano $p\Vm$} Con $T=p(\Vm-b)/A$ y absorbiendo $A$ en la constante:
\[
  p(\Vm-b)\,(\Vm-b)^{A/C}=\cte\quad\Longrightarrow\quad\boxed{p\,(\Vm-b)^{\frac{C+A}{C}}=\cte}.
\]

\paso{Comprobación} Con $b=0$, $A=R$ y $C=\bar C_V$ el exponente es $(\bar C_V+R)/\bar C_V=\gamma$: se recupera
$p\Vm^{\gamma}=\cte$.
\end{resolucion}

% ---------------------------------------------------------------------
\subsection{Problema 27. Expansión isoterma por etapas}
% ---------------------------------------------------------------------
\begin{problema}{27}{Del proceso irreversible al reversible}
Determínese el trabajo en la expansión isoterma de \SI{1}{mol} de gas ideal a \SI{300}{K} de \SI{10}{bar} a
\SI{1}{bar}: (a) en una etapa contra $p_{\mathrm{ext}}=\SI{1}{bar}$; (b) en dos etapas, $10\to5$ contra \SI{5}{bar} y
$5\to1$ contra \SI{1}{bar}; (c) en tres etapas, $10\to5\to2\to1$ contra 5, 2 y \SI{1}{bar}; (d) en infinitas etapas,
con $p_{\mathrm{ext}}$ ligeramente inferior a la presión del gas. Compárense los resultados.
\end{problema}

\begin{resolucion}
\idea Método M2 por etapas: en cada una, la presión exterior es la presión \emph{final} de la etapa y los
volúmenes son los de equilibrio a \SI{300}{K}. El límite de infinitas etapas es el proceso reversible (M1).

\paso{Paso 0: datos} $RT=8{,}314\cdot300=\SI{2494,2}{J}$. Volúmenes de equilibrio $V=RT/p$: 2,494, 4,988,
12,471 y \SI{24,942}{L} a 10, 5, 2 y \SI{1}{bar}.

\paso{Paso 1: una etapa genérica} De $p_i$ a $p_f$ contra $p_{\mathrm{ext}}=p_f$:
\[
  W=-p_f\bigl(V_f-V_i\bigr)=-p_f\Bigl(\frac{RT}{p_f}-\frac{RT}{p_i}\Bigr)=-RT\Bigl(1-\frac{p_f}{p_i}\Bigr).
\]

\paso{Paso 2: los cuatro casos}
\begin{align*}
  W_a&=-RT\Bigl(1-\tfrac{1}{10}\Bigr)=-0{,}9\,RT=\boxed{-\SI{2,24}{kJ}},\\
  W_b&=-RT\Bigl[\Bigl(1-\tfrac{5}{10}\Bigr)+\Bigl(1-\tfrac15\Bigr)\Bigr]=-1{,}3\,RT=\boxed{-\SI{3,24}{kJ}},\\
  W_c&=-RT\Bigl[\Bigl(1-\tfrac{5}{10}\Bigr)+\Bigl(1-\tfrac25\Bigr)+\Bigl(1-\tfrac12\Bigr)\Bigr]=-1{,}6\,RT=\boxed{-\SI{3,99}{kJ}},\\
  W_d&=-RT\ln\frac{p_1}{p_2}=-RT\ln10=-2{,}303\,RT=\boxed{-\SI{5,74}{kJ}}.
\end{align*}

\paso{Paso 3: calores} En los cuatro casos el estado final es el mismo y $T$ no cambia, luego $\Delta U=0$ y
$Q=-W$: el gas absorbe del foco \num{2,24}, \num{3,24}, \num{3,99} y \SI{5,74}{kJ}.

\paso{Comparación} $|W_a|<|W_b|<|W_c|<|W_d|$. Cada etapa añadida acerca $p_{\mathrm{ext}}$ a la presión del gas y
aumenta el área de los rectángulos bajo la isoterma (figura~\ref{fig:p27}); el trabajo entregado es máximo en el
proceso reversible.
\end{resolucion}

\begin{figure}[htbp]
\centering
\begin{tikzpicture}
\begin{groupplot}[group style={group size=2 by 2,horizontal sep=1.3cm,vertical sep=1.9cm},
    graf,width=7.2cm,height=5.0cm,xmin=0,xmax=26.5,ymin=0,ymax=11.2,
    xtick={0,5,10,15,20,25},ytick={1,2,5,10},
    grid=major,grid style={regla!55,line width=0.3pt},
    xlabel={$V$ (L)},ylabel={$p$ (bar)},
    title style={font=\small,align=center,yshift=-3pt}]
\nextgroupplot[title={(a) Una etapa contra \SI{1}{bar}}]
  \addplot[draw=navy,line width=0.9pt,fill=fondomedio] coordinates {(2.494,0) (2.494,1) (24.942,1) (24.942,0)} \closedcycle;
  \addplot[acento,densely dashed,domain=2.494:24.942,samples=90] {24.942/x};
  \addplot[only marks,mark=*,mark size=1.8pt,rojo] coordinates {(2.494,10) (24.942,1)};
  \node[font=\scriptsize,navy,fill=white,inner sep=1.5pt,anchor=south] at (axis cs:16,6.6) {$|W|=0{,}9\,RT=\SI{2,24}{kJ}$};
\nextgroupplot[title={(b) Dos etapas: 5 y \SI{1}{bar}}]
  \addplot[draw=navy,line width=0.9pt,fill=fondomedio] coordinates {(2.494,0) (2.494,5) (4.988,5) (4.988,1) (24.942,1) (24.942,0)} \closedcycle;
  \addplot[acento,densely dashed,domain=2.494:24.942,samples=90] {24.942/x};
  \addplot[only marks,mark=*,mark size=1.8pt,rojo] coordinates {(2.494,10) (4.988,5) (24.942,1)};
  \node[font=\scriptsize,navy,fill=white,inner sep=1.5pt,anchor=south] at (axis cs:16,6.6) {$|W|=1{,}3\,RT=\SI{3,24}{kJ}$};
\nextgroupplot[title={(c) Tres etapas: 5, 2 y \SI{1}{bar}}]
  \addplot[draw=navy,line width=0.9pt,fill=fondomedio] coordinates {(2.494,0) (2.494,5) (4.988,5) (4.988,2) (12.471,2) (12.471,1) (24.942,1) (24.942,0)} \closedcycle;
  \addplot[acento,densely dashed,domain=2.494:24.942,samples=90] {24.942/x};
  \addplot[only marks,mark=*,mark size=1.8pt,rojo] coordinates {(2.494,10) (4.988,5) (12.471,2) (24.942,1)};
  \node[font=\scriptsize,navy,fill=white,inner sep=1.5pt,anchor=south] at (axis cs:16,6.6) {$|W|=1{,}6\,RT=\SI{3,99}{kJ}$};
\nextgroupplot[title={(d) Infinitas etapas (reversible)}]
  \addplot[draw=none,fill=fondomedio,domain=2.494:24.942,samples=90] {24.942/x} \closedcycle;
  \addplot[navy,line width=0.9pt,domain=2.494:24.942,samples=90] {24.942/x};
  \addplot[only marks,mark=*,mark size=1.8pt,rojo] coordinates {(2.494,10) (24.942,1)};
  \node[font=\scriptsize,navy,fill=white,inner sep=1.5pt,anchor=south] at (axis cs:16,6.6) {$|W|=RT\ln10=\SI{5,74}{kJ}$};
\end{groupplot}
\end{tikzpicture}
\caption{Problema 27. En cada etapa el gas se expande contra la presión exterior fija (escalón en azul oscuro) y
acaba en un estado de equilibrio sobre la isoterma de \SI{300}{K} (puntos rojos). El trabajo entregado es el
área rellena bajo $p_{\mathrm{ext}}$; crece con el número de etapas hasta el área bajo la propia isoterma en el
límite reversible (d). La isoterma discontinua solo indica los estados de equilibrio: durante una etapa
irreversible la presión del gas no está definida.}
\label{fig:p27}
\end{figure}

% ---------------------------------------------------------------------
\subsection{Problema 28. Coeficientes calorimétricos y adiabáticas}
% ---------------------------------------------------------------------
\begin{problema}{28}{Gas perfecto diatómico}
(a) Determínense a \SI{1}{bar} y \SI{0}{\celsius} los coeficientes calorimétricos de un mol de gas perfecto diatómico.
(b) Hállense, para un gas perfecto, las ecuaciones de las adiabáticas reversibles en los diagramas $(V,T)$,
$(p,T)$ y $(p,V)$.
\end{problema}

\begin{resolucion}
\idea (a) Método M4 con la tabla~\ref{tab:calorim}; (b) método M5, ecuaciones de Poisson.

\paso{(a) Paso 1: estado} $T=\SI{273,15}{K}$, $p=\SI{e5}{Pa}$,
$\Vm=RT/p=8{,}314\cdot273{,}15/10^{5}=\SI{2,2711e-2}{m^3/mol}$. Diatómico: $\bar C_V=\tfrac52R=\SI{20,79}{J/(mol.K)}$,
$\bar C_p=\tfrac72R=\SI{29,10}{J/(mol.K)}$.

\paso{(a) Paso 2: coeficientes} De la tabla~\ref{tab:calorim} con $n=1$:
\begin{align*}
  l&=p=\boxed{\SI{1,000e5}{Pa}}, &
  h&=-\Vm=\boxed{-\SI{2,271e-2}{m^3}},\\
  \lambda&=\frac{\bar C_V\Vm}{R}=\tfrac52\Vm=\boxed{\SI{5,678e-2}{m^3}}, &
  \mu&=\frac{\bar C_p\,p}{R}=\tfrac72p=\boxed{\SI{3,500e5}{Pa}}.
\end{align*}
Las unidades son las de \eqref{eq:calorim}: $l$ y $\mu$ en \si{J/m^3}$=$\si{Pa}; $h$ y $\lambda$ en \si{J/Pa}$=$\si{m^3}.

\paso{(b) Adiabáticas} Por \eqref{eq:poisson}, con $\gamma=7/5$ para el diatómico:
\[
  (V,T):\ TV^{2/5}=\cte,\qquad (p,T):\ Tp^{-2/7}=\cte,\qquad (p,V):\ pV^{7/5}=\cte,
\]
que son $TV^{\gamma-1}$, $Tp^{(1-\gamma)/\gamma}$ y $pV^{\gamma}$ con $\gamma-1=\tfrac25$ y $(1-\gamma)/\gamma=-\tfrac27$.
La deducción completa está en la sección de adiabáticas reversibles.
\end{resolucion}

% ---------------------------------------------------------------------
\subsection{Problema 29. Expansión adiabática irreversible}
% ---------------------------------------------------------------------
\begin{problema}{29}{Mismos estados, procesos irreversible y reversible}
Dos moles de gas perfecto diatómico, a $p_1=\SI{4}{atm}$ y $T_1=\SI{300}{K}$, se expanden adiabáticamente contra
$p_0=\SI{1}{atm}$ constante hasta triplicar su volumen. Calcúlense $T_2$, $Q$, $W$ y $\Delta U$. Repítase para un
proceso reversible entre los mismos estados formado por una isoterma hasta el volumen final y un isócoro
hasta la presión final.
\end{problema}

\begin{resolucion}
\idea Proceso irreversible: el trabajo se calcula con $p_0$ (M2) y, como $Q=0$, el balance da $\Delta U$ y de él
$T_2$ (M3). Camino reversible: $\Delta U$ es el mismo porque los estados son los mismos; $Q$ y $W$ no.

\paso{Paso 0: estado inicial} $nRT_1=2\cdot8{,}314\cdot300=\SI{4988,4}{J}$ y
$V_1=nRT_1/p_1=4988{,}4/405\,300=\SI{1,2308e-2}{m^3}$; $V_2=3V_1=\SI{3,6924e-2}{m^3}$.

\paso{Proceso irreversible. Trabajo} Contra presión exterior fija,
\[
  W=-p_0(V_2-V_1)=-p_0\cdot2V_1=-2\,\frac{p_0}{p_1}\,nRT_1=-\frac{nRT_1}{2}=\boxed{-\SI{2494}{J}}.
\]

\paso{Proceso irreversible. Balance} $Q=0$ (adiabático), luego $\Delta U=W=-\SI{2494}{J}$. Como
$\Delta U=n\bar C_V\Delta T$ con $\bar C_V=\tfrac52R$:
\[
  \Delta T=\frac{W}{n\cdot\frac52R}=\frac{-nRT_1/2}{\frac52nR}=-\frac{T_1}{5}=-\SI{60}{K}
  \quad\Longrightarrow\quad\boxed{T_2=\SI{240}{K}}.
\]
La presión final es $p_2=nRT_2/V_2=p_1\,\frac{T_2}{T_1}\,\frac{V_1}{V_2}=4\cdot0{,}8\cdot\tfrac13=\SI{1,067}{atm}$.

\paso{Proceso reversible. Isoterma a $T_1$} $\Delta U_1=0$ y
\[
  W_1=-nRT_1\ln\frac{V_2}{V_1}=-4988{,}4\ln3\ \mathrm J=-\SI{5480}{J},\qquad Q_1=-W_1=+\SI{5480}{J}.
\]

\paso{Proceso reversible. Isócoro hasta $p_2$} $W_2=0$ y $Q_2=\Delta U_2=n\bar C_V(T_2-T_1)=-\SI{2494}{J}$.

\paso{Proceso reversible. Totales} $W=-\SI{5480}{J}$, $Q=5480-2494=\SI{2986}{J}$, $\Delta U=Q+W=\boxed{-\SI{2494}{J}}$,
el mismo que en el proceso irreversible, y $T_2=\SI{240}{K}$ por construcción.
\end{resolucion}

\begin{figure}[htbp]
\centering
\begin{tikzpicture}
\begin{axis}[graf,width=9.8cm,height=6.2cm,xmin=6,xmax=56,ymin=0,ymax=4.7,
    xtick={12.308,36.924},xticklabels={$V_1$,$V_2=3V_1$},ytick={1,4},yticklabels={$p_0=1$,4},
    grid=major,grid style={regla!55,line width=0.3pt},
    xlabel={$V$ (L)},ylabel={$p$ (atm)},leyenda abajo,clip=false]
  % trabajo reversible (área bajo la isoterma) y trabajo irreversible (bajo p0)
  \addplot[draw=none,fill=fondomedio,domain=12.308:36.924,samples=80,forget plot] {4*12.308/x} \closedcycle;
  \addplot[draw=none,pattern=north east lines,pattern color=acento,forget plot] coordinates {(12.308,0) (12.308,1) (36.924,1) (36.924,0)} \closedcycle;
  \addplot[navy,line width=1.1pt,domain=12.308:36.924,samples=80] {4*12.308/x};
  \addplot[acento,line width=1.3pt] coordinates {(36.924,1.3333) (36.924,1.0667)};
  \addplot[rojo,densely dashed,line width=1.1pt] coordinates {(12.308,1) (36.924,1)};
  % sentido de recorrido
  \draw[navy,-{Stealth[length=5pt]}] (axis cs:19.0,{4*12.308/19.0+0.22}) -- (axis cs:22.0,{4*12.308/22.0+0.19});
  % estados
  \fill[black] (axis cs:12.308,4) circle (1.8pt);
  \fill[black] (axis cs:36.924,1.3333) circle (1.5pt);
  \fill[black] (axis cs:36.924,1.0667) circle (1.8pt);
  \node[font=\scriptsize,anchor=west,inner sep=3pt] at (axis cs:12.308,4) {$1$: \SI{4}{atm}, \SI{300}{K}};
  \node[font=\scriptsize,anchor=south west,inner sep=2pt] at (axis cs:36.924,1.3333) {$\tfrac43$ atm, \SI{300}{K}};
  \node[font=\scriptsize,anchor=north west,inner sep=2pt] at (axis cs:36.924,1.0667) {$2$: \SI{1,07}{atm}, \SI{240}{K}};
  % áreas
  \node[font=\scriptsize,navy,fill=fondomedio,inner sep=1.5pt] at (axis cs:22.5,1.55) {$|W_{\mathrm{rev}}|=\SI{5,48}{kJ}$};
  \node[font=\scriptsize,navy,fill=white,inner sep=1.5pt] at (axis cs:24.6,0.5) {$|W_{\mathrm{irr}}|=\SI{2,49}{kJ}$};
  \legend{isoterma reversible a \SI{300}{K},isócoro reversible,$p_{\mathrm{ext}}=p_0$ (proceso irreversible)}
\end{axis}
\end{tikzpicture}
\caption{Problema 29. El proceso irreversible no tiene trayectoria en el plano $p$--$V$: solo están definidos sus
estados extremos 1 y 2, y su trabajo es el área bajo $p_0$ (rayada). El camino reversible entre los mismos estados,
isoterma a \SI{300}{K} y después isócoro, entrega el área bajo la isoterma (rellena), mayor, porque recibe calor del
foco; el isócoro no aporta trabajo. $\Delta U=-\SI{2494}{J}$ en los dos casos.}
\label{fig:p29}
\end{figure}

\begin{observacion}[Presión final]
$p_2=\SI{1,067}{atm}>p_0$: el gas no ha alcanzado el equilibrio mecánico con el exterior. El enunciado fija el
volumen final ($3V_1$), lo que exige unos topes que detengan el pistón.
\end{observacion}

% ---------------------------------------------------------------------
\subsection{Problema 30. Capacidad calorífica de un proceso}
% ---------------------------------------------------------------------
\begin{problema}{30}{Capacidad calorífica a lo largo de $\Pi$}
Demuéstrese para un sistema hidrostático
\[
  \frac{\bar C_\Pi}{\bar C_V}=\frac{\kappa_T-\gamma\kappa_\Pi}{\kappa_T-\kappa_\Pi},
\]
donde $\bar C_\Pi$ es el calor molar a lo largo de un proceso $\Pi$, $\gamma=\bar C_p/\bar C_V$ y
$\kappa_\Pi\equiv-\frac1{\Vm}\pdi{\Vm}{p}{\Pi}$. Aplíquese al proceso $\Pi:\ p=b+m\Vm$ ($b>0$, $m<0$) de un gas
perfecto monoatómico.
\end{problema}

\begin{resolucion}
\idea Método M4: escribir $\dbar Q=\bar C_V\dd T+l\,\dd\Vm$ y dividir por $\dd T$ \emph{a lo largo de $\Pi$}; la
derivada $\pdi{\Vm}{T}{\Pi}$ se obtiene combinando la ecuación diferencial de estado con la definición de
$\kappa_\Pi$.

\paso{Paso 1: capacidad del proceso} Por definición, $\bar C_\Pi=(\dbar Q/\dd T)_\Pi$, y con la primera forma de
\eqref{eq:calorim},
\[
  \bar C_\Pi=\bar C_V+l\,\pd{\Vm}{T}{\Pi}.
\]

\paso{Paso 2: el coeficiente $l$} De \eqref{eq:calorimrel}, $\bar C_p=\bar C_V+l\alpha\Vm$, luego
$l=\dfrac{\bar C_p-\bar C_V}{\alpha\Vm}$.

\paso{Paso 3: derivada a lo largo de $\Pi$} En cualquier cambio de estado,
$\dd\Vm=\alpha\Vm\,\dd T-\kappa_T\Vm\,\dd p$ (Tema 1). A lo largo de $\Pi$, además,
$\dd\Vm=-\kappa_\Pi\Vm\,\dd p$, es decir $\dd p=-\dd\Vm/(\kappa_\Pi\Vm)$. Sustituyendo:
\[
  \dd\Vm=\alpha\Vm\,\dd T+\frac{\kappa_T}{\kappa_\Pi}\,\dd\Vm
  \quad\Longrightarrow\quad
  \dd\Vm\Bigl(1-\frac{\kappa_T}{\kappa_\Pi}\Bigr)=\alpha\Vm\,\dd T
  \quad\Longrightarrow\quad
  \pd{\Vm}{T}{\Pi}=\frac{\alpha\Vm\,\kappa_\Pi}{\kappa_\Pi-\kappa_T}.
\]

\paso{Paso 4: resultado} Llevando los pasos 2 y 3 al 1, el factor $\alpha\Vm$ se cancela:
\[
  \bar C_\Pi=\bar C_V+(\bar C_p-\bar C_V)\frac{\kappa_\Pi}{\kappa_\Pi-\kappa_T}
  \;\Longrightarrow\;
  \frac{\bar C_\Pi}{\bar C_V}=1+(\gamma-1)\frac{\kappa_\Pi}{\kappa_\Pi-\kappa_T}
  =\frac{\gamma\kappa_\Pi-\kappa_T}{\kappa_\Pi-\kappa_T}
  =\boxed{\frac{\kappa_T-\gamma\kappa_\Pi}{\kappa_T-\kappa_\Pi}}.
\]
Comprobaciones: isóbaro, $\kappa_\Pi\to\infty$ y $\bar C_\Pi\to\gamma\bar C_V=\bar C_p$; isócoro, $\kappa_\Pi=0$ y
$\bar C_\Pi=\bar C_V$; isotermo, $\kappa_\Pi=\kappa_T$ y $\bar C_\Pi\to\infty$.

\paso{Paso 5: aplicación} Gas perfecto monoatómico: $\gamma=5/3$, $\bar C_V=\tfrac32R$, $\kappa_T=1/p$. Para
$p=b+m\Vm$, $\pdi{\Vm}{p}{\Pi}=1/m$ y $\kappa_\Pi=-\dfrac{1}{m\Vm}$. Multiplicando numerador y denominador por
$pm\Vm$:
\[
  \frac{\bar C_\Pi}{\bar C_V}=\frac{\frac1p+\frac{5}{3m\Vm}}{\frac1p+\frac{1}{m\Vm}}
  =\frac{m\Vm+\frac53p}{m\Vm+p}
  \;\xrightarrow{\;p=b+m\Vm\;}\;
  \frac{\frac83m\Vm+\frac53b}{2m\Vm+b}
  \quad\Longrightarrow\quad
  \boxed{\bar C_\Pi=\frac R2\,\frac{8m\Vm+5b}{2m\Vm+b}}.
\]
\end{resolucion}

\begin{figure}[htbp]
\centering
\begin{tikzpicture}
\begin{axis}[graf,width=8.0cm,height=5.6cm,xmin=0,xmax=1,ymin=-6,ymax=8,
    xtick={0,0.5,0.625,1},xticklabels={0,$\tfrac12$,$\tfrac58$,1},ytick={-4,0,4,8},
    xlabel={$u=-m\Vm/b$},ylabel={$\bar C_\Pi/R$}]
  \fill[fondomedio] (axis cs:0.5,-6) rectangle (axis cs:0.625,8);
  \draw[gris,densely dotted] (axis cs:0.5,-6) -- (axis cs:0.5,8);
  \draw[gris,line width=0.4pt] (axis cs:0,0) -- (axis cs:1,0);
  \addplot[navy,domain=0:0.478,samples=80] {(5-8*x)/(2*(1-2*x))};
  \addplot[navy,domain=0.522:1,samples=80] {(5-8*x)/(2*(1-2*x))};
  \fill[rojo] (axis cs:0.625,0) circle (1.6pt);
  \fill[navy] (axis cs:0,2.5) circle (1.6pt) node[anchor=north west,font=\scriptsize,inner sep=2pt] {$\bar C_p$};
  \fill[navy] (axis cs:1,1.5) circle (1.6pt) node[anchor=south east,font=\scriptsize,inner sep=3pt] {$\bar C_V$};
  \node[font=\scriptsize,anchor=south west,align=left] at (axis cs:0.04,5.2) {absorbe calor,\\$T$ sube};
  \draw[gris,line width=0.4pt] (axis cs:0.6,-4.7) -- (axis cs:0.67,-4.2);
  \node[font=\scriptsize,anchor=west,align=left,inner sep=2pt] at (axis cs:0.67,-4.0) {absorbe calor,\\$T$ baja};
  \node[font=\scriptsize,anchor=south,align=center] at (axis cs:0.8,4.4) {cede calor,\\$T$ baja};
\end{axis}
\end{tikzpicture}
\caption{Problema 30. Capacidad molar a lo largo de la recta $p=b+m\Vm$ en función de $u=-m\Vm/b\in(0,1)$.
En $u=\tfrac12$ la temperatura es máxima ($T\propto p\Vm$) y $\bar C_\Pi$ diverge; en la franja sombreada
$\bar C_\Pi<0$; en $u=\tfrac58$ la recta es tangente a
una adiabática y $\bar C_\Pi=0$. Para $u\to0$ el proceso es casi isóbaro y $\bar C_\Pi\to\bar C_p=\tfrac52R$.}
\label{fig:p30}
\end{figure}

\begin{observacion}[Lectura del signo]
Con $\dbar Q=\bar C_\Pi\,\dd T$: para $u<\tfrac12$, $\dd T>0$ y $\bar C_\Pi>0$, el gas absorbe calor; entre
$\tfrac12$ y $\tfrac58$, $\dd T<0$ y $\bar C_\Pi<0$, sigue absorbiendo aunque se enfríe; a partir de
$u=\tfrac58$ cede calor. Una capacidad calorífica negativa no es un error: describe un proceso en el que el gas
entrega más trabajo del calor que recibe.
\end{observacion}

% ---------------------------------------------------------------------
\subsection{Parcial 2025/26, cuestión 2}
% ---------------------------------------------------------------------
\begin{problema}{P.2}{Coeficiente adiabático de un gas no ideal}
Un gas tiene $pV=AT+B$ y $U=CT+DV$, con $A$, $B$, $C$ y $D$ constantes positivas. (a) Calcule $\gamma\equiv C_p/C_V$.
(b) Muestre que para presiones grandes tiene un valor constante mayor que 1.
\end{problema}

\begin{resolucion}
\idea Método M4: $C_V$ de $U$; $C_p=C_V+l\,\pdi VTp$ con $l=\pdi UVT+p$.

\paso{(a) Paso 1: $C_V$} $C_V=\pdi UTV=C$.

\paso{(a) Paso 2: $l$ y la derivada} $\pdi UVT=D$, luego $l=D+p$. De $V=(AT+B)/p$, a $p$ fija,
$\pdi VTp=A/p$.

\paso{(a) Paso 3: $C_p$ y $\gamma$}
\[
  C_p=C+(D+p)\frac Ap=C+A+\frac{AD}p
  \quad\Longrightarrow\quad
  \boxed{\gamma=\frac{C_p}{C_V}=1+\frac AC+\frac{AD}{Cp}=\frac{(C+A)\,p+AD}{C\,p}}.
\]

\paso{(b) Límite} Para $p\to\infty$ el último término se anula:
\[
  \gamma\;\longrightarrow\;1+\frac AC=\frac{C+A}{C},
\]
constante y mayor que 1 porque $A,C>0$. Es el mismo exponente que en el problema 26.
\end{resolucion}

% ---------------------------------------------------------------------
\subsection{Parcial 2025/26, cuestiones 3(a) y 4(a)}
% ---------------------------------------------------------------------
\begin{problema}{P.3(a)}{Expansión isoterma: la parte de la primera ley}
Un mol de gas ideal se expande isoterma y reversiblemente desde $(p_i,V_i)$ hasta $p_i/2$, en contacto con un
foco a $T_i$. ¿Cuál es su cambio de energía? (El cambio de entropía se trata en el Tema 3.)
\end{problema}

\begin{resolucion}
\idea Gas ideal isotermo: $\Delta U=0$ y $Q=-W$ (tabla~\ref{tab:procesos}).

\paso{Paso 1} Por la ley de Joule, $\boxed{\Delta U=0}$. Con $p_iV_i=p_fV_f$, $V_f=2V_i$.

\paso{Paso 2} $W=-RT_i\ln(V_f/V_i)=-RT_i\ln2$ y $Q=-W=RT_i\ln2>0$: el gas absorbe del foco exactamente el
trabajo que entrega.
\end{resolucion}

\begin{problema}{P.4(a)}{Calor en el isócoro de un ciclo}
Un gas perfecto monoatómico se calienta adiabática y reversiblemente de $T_1$ a $T_2$ reduciendo su volumen a la
mitad ($1\to2$); se expande en contacto con un foco a $T_2$ hasta su volumen inicial ($2\to3$); y vuelve al estado
inicial en contacto con un foco a $T_1$ mediante un proceso irreversible a volumen constante ($3\to1$). ¿Cuál es el
calor intercambiado en $3\to1$?
\end{problema}

\begin{resolucion}
\idea Aunque $3\to1$ sea irreversible, es isócoro: $W=0$ y $Q=\Delta U$, que solo depende de los estados.

\paso{Paso 1: temperaturas} $T_3=T_2$ (isoterma) y, por Poisson con $\gamma-1=2/3$ y $V_1/V_2=2$,
$T_2=T_1\cdot2^{2/3}=1{,}587\,T_1$.

\paso{Paso 2: balance del isócoro} Con $W_{31}=0$,
\[
  Q_{31}=\Delta U_{31}=n\bar C_V(T_1-T_3)=\boxed{\tfrac32nR\,(T_1-T_2)}=\tfrac32nRT_1\bigl(1-2^{2/3}\bigr)\simeq-0{,}881\,nRT_1 .
\]
Es negativo: el gas cede calor al foco frío. El cambio de entropía del universo y el teorema de Carnot,
apartados (b) y (c), se tratan en el Tema 3.
\end{resolucion}

% ---------------------------------------------------------------------
\subsection{Resumen de resultados}
% ---------------------------------------------------------------------
\begin{table}[H]
\centering
\caption{Resultados de los seminarios del Tema 2.}
\small
\renewcommand{\arraystretch}{1.3}
\begin{tabularx}{\textwidth}{lXc}
\toprule
Problema & Resultado & Método\\
\midrule
17 & $a=\SI{2,375}{MPa}$, $b=-\SI{75}{MPa/m^3}$;\enspace $W=+\SI{25}{kJ}$ & M1\\
18 & $W=-30$; $-20{,}79$; $-\SI{17,57}{kJ}$ para $a=0$; 1; 1,5 & M1\\
19 & $-nRT\ln\frac{2V_0-nb}{V_0-nb}$;\enspace $-nRT\bigl(\ln2+\frac{nB}{2V_0}\bigr)$;\enspace $-nRT\ln\frac{2V_0-nb}{V_0-nb}+\frac{n^2a}{2V_0}$ & M1\\
20 & $T_2=\SI{149,1}{K}$;\enspace $n=\SI{1,793}{mol}$ & M5, M3\\
21 & $W=-\SI{15}{kJ}$;\enspace $\Delta U=-\SI{10}{kJ}$ & M1, M3\\
22 & $W_{\mathrm{conf}}=\SI{900}{J}$;\enspace $\Delta z=-\SI{0,259}{m}$ (baja) & M2, M3\\
23 & $W=-Q$: $2RT\ln3$;\enspace $\tfrac83RT$;\enspace $4RT$ & M1, M3\\
24 & $Q=2aT+3bT^2-c/T$ & M4\\
25 & $\bar C_V=c$;\enspace $\bar C_p=c+\dfrac{R^2T\Vm^3}{RT\Vm^3-2a(\Vm-b)^2}$ & M4\\
26 & $p(\Vm-b)^{(C+A)/C}=\cte$ & M5\\
27 & $W=-0{,}9RT$; $-1{,}3RT$; $-1{,}6RT$; $-RT\ln10$ ($-2{,}24$; $-3{,}24$; $-3{,}99$; $-\SI{5,74}{kJ}$) & M2, M1\\
28 & $l=\SI{e5}{Pa}$, $h=-\SI{2,27e-2}{m^3}$, $\lambda=\SI{5,68e-2}{m^3}$, $\mu=\SI{3,5e5}{Pa}$;\enspace $TV^{2/5}$, $Tp^{-2/7}$, $pV^{7/5}$ & M4, M5\\
29 & Irrev.: $T_2=\SI{240}{K}$, $Q=0$, $W=\Delta U=-\SI{2494}{J}$.\enspace Rev.: $W=-\SI{5480}{J}$, $Q=\SI{2986}{J}$ & M2, M3\\
30 & $\bar C_\Pi/\bar C_V=(\kappa_T-\gamma\kappa_\Pi)/(\kappa_T-\kappa_\Pi)$;\enspace $\bar C_\Pi=\frac R2\frac{8m\Vm+5b}{2m\Vm+b}$ & M4\\
P.2 & $\gamma=1+A/C+AD/(Cp)\to1+A/C>1$ & M4\\
P.3(a) & $\Delta U=0$;\enspace $Q=-W=RT_i\ln2$ & M3\\
P.4(a) & $Q_{31}=\tfrac32nR(T_1-T_2)\simeq-0{,}881\,nRT_1$ & M3, M5\\
\bottomrule
\end{tabularx}
\end{table}

\begin{recuadro}{Lista de control antes de entregar}
\begin{itemize}[label={\color{acento}$\blacktriangleright$},leftmargin=1.4em,itemsep=4pt]
\item ¿Signos con la convención del curso? Expansión, $W<0$; calor cedido, $Q<0$.
\item ¿El proceso es reversible? Si no, el trabajo va con $p_{\mathrm{ext}}$, nunca con la presión del gas.
\item ¿Unidades SI (\si{Pa}, \si{m^3}, \si{K})? $\SI{1}{L}=\SI{e-3}{m^3}$, $\SI{1}{bar}=\SI{e5}{Pa}$; en $\Delta T$, °C y K coinciden.
\item ¿Se ha usado $\Delta U=nC_V\Delta T$ en un gas perfecto aunque el proceso no sea isócoro? Es correcto.
\item En calorimetría, ¿se ha incluido $\pdi UVT$ en $l$ si el gas no es ideal?
\item ¿Se ha comprobado un límite (gas ideal, Mayer, $\gamma$) o el mismo $\Delta U$ por dos caminos?
\end{itemize}
\end{recuadro}

